% =================================================
% Теоретический лист. Урок 3
% =================================================

\worksheettitle
  {Составление чисел. Задачи на цифры}
  {Урок 3 из 3. Обозначение натуральных чисел}

\begin{problembox}[Главный вопрос]
  Как составить \textbf{все} числа, удовлетворяющие условию, ничего не
  пропустить и не записать лишнего?
\end{problembox}

В таких заданиях важно не угадывать ответы, а организовать
\textbf{упорядоченный перебор}: выбрать первую цифру, затем последовательно
рассмотреть все допустимые продолжения.

\section{Порядок цифр имеет значение}

Одни и те же цифры могут образовать разные числа.

\orderchangesfigure

\begin{propertybox}[Позиционный принцип]
  В числе значение цифры зависит от её места. Поэтому перестановка цифр
  обычно изменяет число.
\end{propertybox}

Например, из цифр \(2\) и \(7\) можно получить числа \(27\) и \(72\).
В первом числе \(2\) обозначает два десятка, а во втором --- две единицы.

\section{Когда цифры могут повторяться}

Составим все двузначные числа, в записи которых используются только цифры
\(3\) и \(7\). Повторение цифр разрешено.

\repeatstreefigure

Сначала выбираем цифру десятков:
\[
  3\quad\text{или}\quad 7.
\]
Для каждого выбора снова рассматриваем обе возможные цифры единиц:
\[
  \boxed{33},\quad \boxed{37},\quad \boxed{73},\quad \boxed{77}.
\]

\begin{conclusionbox}[Как не пропустить число]
  Закончите все варианты для одной цифры десятков и только затем переходите
  к следующей цифре десятков.
\end{conclusionbox}

Если на каждое из двух мест можно поставить одну из двух цифр, получаем
\[
  2\cdot2=4
\]
варианта. На этом этапе достаточно понимать произведение как запись:
«два выбора для десятков и для каждого из них два выбора для единиц».

\section{Когда цифры не должны повторяться}

Составим все трёхзначные числа из цифр \(0\), \(5\), \(8\), используя
каждую цифру ровно один раз.

\norepeattreefigure

Ноль не может быть первой цифрой трёхзначного числа. Поэтому первая цифра
равна \(5\) или \(8\). После её выбора две оставшиеся цифры можно поставить
в двух порядках:
\[
  \boxed{508},\quad \boxed{580},\quad \boxed{805},\quad \boxed{850}.
\]

\begin{warningbox}[Запись не может начинаться с нуля]
  Запись \(053\) обозначает число \(53\), а не трёхзначное число. Начальный
  ноль не создаёт нового разряда.
\end{warningbox}

\section{Наибольшее и наименьшее число}

Чтобы получить наибольшее число, цифры располагают по убыванию.
Чтобы получить наименьшее число, цифры располагают по возрастанию, но ноль
нельзя ставить первым.

\largestsmallestfigure

Для цифр \(0,2,5,8\):
\[
  \text{наибольшее число }8520,
  \qquad
  \text{наименьшее число }2058.
\]

\begin{propertybox}[Правило для наименьшего числа с нулём]
  На первое место поставьте наименьшую ненулевую цифру. Ноль поставьте сразу
  после неё. Остальные цифры расположите по возрастанию.
\end{propertybox}

\section{Условия, связывающие цифры}

Найдём все двузначные числа, в которых цифра десятков на \(3\) больше цифры
единиц.

Обозначать цифры буквами необязательно. Можно начать с цифры единиц и
заполнить таблицу.

\begin{center}
\renewcommand{\arraystretch}{1.35}
\begin{tabular}{c|ccccccc}
  \toprule
  цифра единиц & 0 & 1 & 2 & 3 & 4 & 5 & 6 \\
  \midrule
  цифра десятков & 3 & 4 & 5 & 6 & 7 & 8 & 9 \\
  \bottomrule
\end{tabular}
\end{center}

Получаем:
\[
  \boxed{30,\ 41,\ 52,\ 63,\ 74,\ 85,\ 96}.
\]
После цифры единиц \(6\) продолжать нельзя: к \(7\) пришлось бы прибавить
\(3\), а цифры \(10\) не существует.

\systematicroutefigure

\begin{conclusionbox}[Алгоритм перебора]
  \begin{enumerate}
    \item Уточните, какие цифры и сколько раз можно использовать.
    \item Выберите цифру, которую удобно менять по порядку: \(0,1,2,\ldots\).
    \item По условию найдите остальные цифры.
    \item Исключите невозможные записи: ведущий ноль, повторение или цифру
      больше \(9\).
    \item Проверьте каждое число и убедитесь, что перебор закончен.
  \end{enumerate}
\end{conclusionbox}

\section{Сколько существует подходящих чисел}

Рассмотрим вопрос: сколько двузначных чисел имеют сумму цифр, равную \(7\)?

Цифра десятков не может равняться нулю. Последовательно выбираем её:
\[
\begin{array}{c|ccccccc}
  \text{десятки} & 1&2&3&4&5&6&7\\
  \hline
  \text{единицы} & 6&5&4&3&2&1&0
\end{array}
\]

Получаем числа
\[
  16,\ 25,\ 34,\ 43,\ 52,\ 61,\ 70.
\]
Их \(7\).

\placevaluecountfigure

\begin{notebox}[Список нужен не всегда]
  Если вопрос требует только количество, иногда достаточно посчитать
  допустимые варианты. Но в начале изучения темы полезно выписывать числа:
  так легче проверить, что ни один вариант не потерян.
\end{notebox}

\section{Смешанные задачи на разрядный состав}

Иногда число задают не отдельными цифрами, а количеством десятков, сотен
или тысяч. Тогда сначала переводят каждую часть в обычные единицы.

Например:
\[
  5\text{ десятков тысяч}+78\text{ сотен}
  =5\cdot10\,000+78\cdot100
\]
\[
  =50\,000+7\,800=57\,800.
\]

\begin{warningbox}[Складываем значения, а не приписываем записи]
  Нельзя просто приписать \(78\) к \(5\). Выражения «десяток тысяч» и
  «сотня» обозначают единицы счёта разного размера.
\end{warningbox}

\begin{practicebox}[Проверьте себя]
  \begin{enumerate}
    \item Запишите все двузначные числа из цифр \(2\) и \(6\), если цифры
      могут повторяться.
    \item Запишите все трёхзначные числа из цифр \(0,4,9\), если цифры не
      повторяются и все три цифры используются.
    \item Составьте наибольшее и наименьшее четырёхзначные числа из цифр
      \(0,3,5,8\) без повторений.
    \item Запишите все двузначные числа, в которых цифра десятков на \(2\)
      больше цифры единиц.
  \end{enumerate}
\end{practicebox}

\begin{notebox}[Ответы]
  1) \(22,26,62,66\);\quad
  2) \(409,490,904,940\);\quad
  3) \(8530\) и \(3058\);\quad
  4) \(20,31,42,53,64,75,86,97\).
\end{notebox}

\begin{questionsbox}
  \begin{enumerate}
    \item Почему запись \(053\) не является трёхзначным числом?
    \item Чем отличаются задачи «цифры могут повторяться» и «цифры не
      повторяются»?
    \item Как убедиться, что выписаны все подходящие числа?
    \item Как составить наименьшее число, если среди данных цифр есть ноль?
  \end{enumerate}
\end{questionsbox}
